\chapter{Zero-Copy Inter-Process Communication}
\label{ch:zero_copy_ipc}

The AMPS engine is logically fractured: Python's \texttt{asyncio} event loop orchestrates the 400 generative agents on one set of CPU cores, while the Numba JIT engine computes the Limit Order Book on another. To bridge these two distinct environments, the simulation requires high-throughput Inter-Process Communication (IPC). Traditional IPC mechanisms---such as WebSockets, gRPC, or Redis---introduce serialization and deserialization overhead that violates our microsecond latency budgets. This chapter details the implementation of true zero-copy IPC using shared memory and Apache Arrow.

\section{The Overhead of Serialization}

When two isolated processes communicate in standard Python, the data must be serialized into a byte stream (e.g., using \texttt{pickle} or JSON), transmitted over a local socket, and then deserialized back into Python objects by the receiving process. 

For 400 agents requesting the complete LOB state matrix 1,000 times per second, this serialization/deserialization cycle consumes 99\% of the CPU time, causing the simulation to grind to a halt.

\section{Pure Python Shared Memory (\texttt{/dev/shm})}

To eliminate serialization, AMPS employs Python 3.8+ \texttt{multiprocessing.shared\_memory}. This library maps directly to the POSIX \texttt{/dev/shm} hardware register on Linux, creating a block of physical RAM accessible to multiple isolated processes.

By defining a contiguous memory layout using a \texttt{NumPy dtype}, we achieve a highly structured payload. For example, an agent's financial intent is represented as a 32-byte structure. This structure is mathematically aligned to fit exactly two intents into a single 64-byte L1 cache line on the Gracemont E-Cores, preventing partial cache line fetches.

\begin{figure}[htbp]
\centering
\begin{tikzpicture}[
    scale=0.9, transform shape,
    mem/.style={draw, rectangle, minimum height=1cm, minimum width=2.5cm, align=center},
    arr/.style={->, thick, >=stealth}
]
    \node[mem, fill=blue!10] (agent) {Agent Process\\(Python \texttt{asyncio})};
    
    \node[mem, fill=yellow!20, minimum width=4cm] (shm) [right=2cm of agent] {\texttt{/dev/shm}\\(Physical RAM Block)};
    
    \node[mem, fill=green!10] (engine) [right=2cm of shm] {Matching Engine\\(Numba LLVM)};
    
    \draw[arr] (agent.east) -- node[above] {Writes} node[below] {(No Pickle)} (shm.west);
    \draw[arr] (shm.east) -- node[above] {Reads} node[below] {(Direct Pointers)} (engine.west);
\end{tikzpicture}
\caption{The Zero-Copy IPC pipeline. Data remains physically stationary in \texttt{/dev/shm}; only memory pointers are exchanged between the isolated processes.}
\label{fig:shm_pipeline}
\end{figure}

Because Numba natively understands NumPy \texttt{dtype} configurations, it compiles array accesses directly into LLVM memory pointers. The Python agent writes its intent into the \texttt{shared\_memory} block, and the Numba engine reads it directly via pointers. The data never leaves the shared memory, achieving true zero-copy data transit.

\section{The Apache Arrow C Data Interface}

While \texttt{/dev/shm} works perfectly for simple intents, passing the massive, multi-dimensional LOB state matrix requires a more sophisticated memory format. AMPS leverages the \textbf{Apache Arrow C Data Interface}.

Arrow provides a standardized, columnar memory format that acts as a low-level lingua franca for sharing tabular data at runtime. Instead of relying on the heavy \texttt{PyArrow} library, AMPS utilizes Python's \texttt{PyCapsule} objects. 

The Numba matching engine constructs the LOB state in Arrow's columnar format within \texttt{/dev/shm}. It then exposes the underlying C pointers of the \texttt{ArrowSchema} and \texttt{ArrowArray} structures. The Python agents consume these pointers via \texttt{PyCapsule}, instantly accessing the massive LOB matrix without a single byte of data being copied or serialized.

\section{Lock-Free Ring Buffers and MESI Cache Protocols}

Handling concurrent memory access normally requires Mutex locks to prevent data corruption. However, Mutex locks force the operating system to pause threads, introducing massive latency.

To resolve this, AMPS implements \textbf{Lock-Free Ring Buffers}. A ring buffer is a circular array of shared memory managed by atomic integer counters (a \texttt{head} pointer and a \texttt{tail} pointer). The matching engine pushes LOB updates to the \texttt{head}, and agents read from the \texttt{tail}.

When updating these atomic pointers, ownership of the CPU cache line holding the pointers is critical. If multiple CPU cores try to modify the pointer simultaneously, the CPU's MESI (Modified, Exclusive, Shared, Invalid) cache protocol will force the cache line to continuously bounce back and forth between the L1 caches of the physical cores, destroying throughput (a phenomenon known as \textit{False Sharing}). 

By padding the atomic pointers so they sit on independent 64-byte cache lines, and utilizing strict LLVM memory fencing (\texttt{fence release} and \texttt{fence acquire}), AMPS guarantees lock-free, atomic memory visibility across the heterogeneous CPU cores with zero MESI cache bouncing.
